\begin{frame}{Konstrukcija pravokotnice}
	\begin{columns}
  \begin{column}{0.55\textwidth}
	\begin{itemize}
		\item<1-> Dana sta premica \(p\) in točka \(T\).
		\item<2-> Nariši lok \(k\) s središčem v \(T\).
		\item<3-> Lok \(k\) seče premico \(p\) v točkah \(A\) in \(B\).
		\item<4-> Nariši lok \(m\) s središčem v \(A\).
		\item<5-> Nariši lok \(n\) s središčem v \(B\) in z enakim polmerom.
		\item<6-> Loka se sečeta v točki \(C\).
		\item<7> Premica \(TC\) je pravokotna na \(p\).
	\end{itemize}
	 \end{column}
\begin{column}{0.45\textwidth}
	\includegraphics <1>[width=50mm]{slike/fig-1.png}
	\includegraphics <2>[width=50mm]{slike/fig-2.png}
	\includegraphics <3>[width=50mm]{slike/fig-3.png}
	\includegraphics <4>[width=50mm]{slike/fig-4.png}
	\includegraphics <5>[width=50mm]{slike/fig-5.png}
	\includegraphics <6>[width=50mm]{slike/fig-6.png}
	\includegraphics <7>[width=50mm]{slike/fig-7.png}
 \end{column}
\end{columns}
\end{frame}
